% BOOK-FILE-ID: f-305150da53bc
% chapter: 実際の解き方

\chapter{解法の選択と演習}\label{章：解法の選択と演習}
\begin{bookcolumns}

  \chapterref{章：漸化式の解法}では,パターンごとに解法を学んだ.
  しかし,実際の問題では「これは何型です」と教えてもらえるわけではない.
  漸化式を前にしたとき,何から手を付ければよいだろうか.
  基本は「式全体」「数列の項」「付加項」の3つの観点から目立つ特徴を見つけ,それを消すような変形を試すことである.
  この章では,式の特徴から自分で解法を選ぶ練習をしよう.
  1つの式に複数の解法が使えることもあれば,どの解法もうまくいかずに行き詰まることもある.
  変形して得られた式に続けて別の解法を使うこともある.
  例えば分数型で逆数を取った後に特性方程式型が現れれば,次は不動点を引いて等比型へ進む.
  置き換えた式が等差型であれば,そのまま一般項が書ける.
  1回の変形で既知の形にならなくても,そこで得られた式が次の解法の入口になる.
  そうした経験を積むことこそが,\sectionref{節：目指す形}で述べた勘を鍛えることにつながるのだ.
  \BookImport{f-fe20e61c40b1}
  \BookImport{f-bcb250cd9e78}
  \BookImport{f-d44b8bbd86e8}

  % 問題演習を読み込む。§計算と検算は §演習 の中で読み込む。
  \BookImport{f-8dd4179c8f0e}
\end{bookcolumns}
